Training data was drawn from the SFT split of the Ganit dataset
\parencite{ganitllm}, a Bengali mathematical reasoning corpus of 16,868
problems, each carrying a full Bengali reasoning trace and a difficulty score.
The difficulty score is the number of times Qwen3-32B solved the problem across
32 sampled attempts, so it records difficulty relative to a 32B model. The pool
is assembled from four upstream sources: numina-math-cot-bn (13,160 problems),
SOMADHAN (3,487), mCoT-MATH-bn (193) and s1k-32-Bangla (28). Its own difficulty
tags skew easy (10,065 easy, 6,411 olympiad, 258 hard, 134 medium), but the
solve-rate column tells a harsher story: 5,645 problems, a third of the corpus,
were never solved by Qwen3-32B in any of 32 attempts, and 91.7\% of that
unsolvable band comes from numina-math-cot-bn alone. We treat the zero band as
unreliable rather than merely hard, and exclude it.

Figure~\ref{fig:curation} shows the pipeline. Decontamination ran first, before
any difficulty filtering, against 4,080 problems pooled from three public
Bengali mathematics benchmarks: Bn-MGSM (250 problems), Bn-MSVAMP (1,000) and
BenNumEval (2,830 across five task types). A training problem was removed if it
matched a benchmark problem on either of two signals. The first is a numeric
fingerprint, the hash of the sorted set of numbers appearing in the normalised
problem text. Numbers survive re-translation even when every surrounding word
changes, so this catches exactly the case that n-gram or MinHash deduplication
misses. The second is token Jaccard overlap at a threshold of 0.7, which
catches near-identical wording. Normalisation maps Bengali numerals to Western
digits, lowercases, strips punctuation and collapses whitespace, so the two
numeral conventions present in the corpus compare on equal footing. Together
these removed 3,315 problems and left 13,553.

Following LIMO's selection criterion, we then retained problems in the window
$1 \leq \text{correct\_counts} \leq 16$, which reduced 13,553 to 873. A final
near-duplicate pass within the surviving set removed nothing, which is itself
mild evidence that the pool is not padded with reworded repeats. The resulting
set is heavily weighted towards hard problems: 74.6\% carry an olympiad
difficulty tag and 25.3\% are tagged hard, with a mean pass rate of 3.3 out of
32.

A second training set was constructed later to test whether that difficulty
distribution explains our results. Holding every other choice constant, we
selected 1,000 traces from the band $\text{correct\_counts} \geq 28$, giving a
mean pass rate of 31.8 out of 32 and a difficulty composition of 100\% easy.
Section~\ref{sec:difficulty} reports that experiment. This set was
decontaminated directly against the held-out evaluation items rather than
against the public benchmarks, using the same numeric fingerprint idea. That
filter is deliberately aggressive and removed 1,782 candidates.

The evaluation set was independently checked against the 873 training traces by
identifier match, exact normalised text match, and numeric fingerprint combined
with Jaccard overlap. One item matched and was dropped. Neither training set
shares an identifier or problem text with the evaluation set, and the two
training sets are disjoint from each other because they are drawn from
non-overlapping difficulty bands. One limitation is worth stating plainly:
Ganit carries no upstream English source identifiers, so fingerprint and
Jaccard matching are proxies, and a sufficiently free paraphrase could survive
both filters.

\begin{figure}[htbp]
\centering
\includegraphics[width=\linewidth]{figs/fig_curation.pdf}
\caption{Curation of the 873-trace training set. Decontamination against three
public Bengali mathematics benchmarks precedes difficulty filtering, so the
LIMO window is applied to an already-clean pool.}
\label{fig:curation}
\end{figure}
